Full row rank permits extending the \(p\) rows of \(C\) to a basis \(T_0\in\operatorname{GL}_n(\mathbb R)\); write \(W_0=T_0^{-1}\). The rows of \(C_{-x,:}\) are independent. For a fixed column \(j\), restriction of those rows to all coordinates other than \(j\) loses rank precisely when a nonzero linear combination is supported only on coordinate \(j\). Hence
\[
 \operatorname{rank}(C_{-x,-j})<p-1
 \quad\Longleftrightarrow\quad
 e_j^\top\in\operatorname{rowspan}(C_{-x,:}). \tag{15}
\]
The coefficients of \(e_j^\top\) in the row basis \(T_0\) form the row \(e_j^\top W_0\). Since the rows of \(T_0\) indexed by \(\mathsf O\) are the rows of \(C\), (15) becomes
\[
 \operatorname{rank}(C_{-x,-j})<p-1
 \quad\Longleftrightarrow\quad
 (W_0)_{jx}=0\ \text{ and }\ (W_0)_{j\ell}=0
 \text{ for every }\ell\in\mathsf L. \tag{16}
\]

Now let \(j\in J_x(C)\). If \((W_0)_{jx}\ne0\), put \(T_j=T_0\) and \(W_j=W_0\). Otherwise, (16) and the full deletion rank supply \(\ell\in\mathsf L\) with \((W_0)_{j\ell}\ne0\). Define
\[
 T_j=(I_n+e_\ell e_x^\top)T_0,
 \qquad
 W_j=W_0(I_n-e_\ell e_x^\top). \tag{17}
\]
Because \(\ell\ne x\), the two shear matrices in (17) are inverse. The shear changes only latent row \(\ell\), so \((T_j)_{\mathsf O,:}=C\), and
\[
 (W_j)_{jx}=-(W_0)_{j\ell}\ne0. \tag{18}
\]
In the square case \(n=p\), the latent set is empty and (16) shows directly that the first case must hold.

Admissibility also gives \(C_{xj}\ne0\). Since \(W_j^{-1}=T_j\), the inverse-cofactor identity now gives
\[
 (-1)^{j+x}\det((W_j)_{-j,-x})
 =\operatorname{cof}(W_j)_{jx}
 =\det(W_j)(T_j)_{xj}
 =\det(W_j)C_{xj}\ne0. \tag{19}
\]
At least one term in the determinant expansion of the complementary matrix in (19) is nonzero. Its positions form a perfect matching between all rows other than \(j\) and all columns other than \(x\). Augment this matching by assigning row \(j\) to column \(x\), and let \(\Pi_j\) permute rows according to this assignment. Every diagonal entry of \(\Pi_jW_j\) is nonzero. Therefore
\[
 D_j=\operatorname{diag}(\Pi_jW_j),
 \qquad
 Q_j=D_j^{-1}\Pi_jW_j,
 \qquad
 H_j=D_j^{-1}\Pi_j \tag{20}
\]
are well defined. The matrix \(Q_j\) is nonsingular and has unit diagonal, while \(H_j\) is monomial and assigns source \(j\) to equation \(x\). Direct multiplication gives
\[
 Q_j^{-1}H_j=W_j^{-1}\Pi_j^{-1}D_jD_j^{-1}\Pi_j=T_j,
 \qquad
 [Q_j^{-1}H_j]_{\mathsf O,:}=C. \tag{21}
\]
Thus the observed vector generated from \(\varepsilon\) has law \(P_X\). The anchor source assumptions and the nonzero, pairwise nonproportional columns of \(C\) show that \((n,C,\varepsilon)\in\mathfrak I(P_X)\), so (20)--(21) meet the observational-law and irreducibility clauses of \(\mathfrak F_x(P_X)\).

It remains to check the intervention boundary rather than infer it from observational solvability. Since row \(j\) was assigned to equation \(x\),
\[
 (Q_j^{-1})_{\mathsf O,x}
 =(T_j)_{\mathsf O,j}(D_j)_{xx}
 =C_{:,j}(D_j)_{xx}. \tag{22}
\]
Both \(C_{xj}\) and \((D_j)_{xx}\) are nonzero, so \((Q_j^{-1})_{xx}\ne0\). The principal cofactor identity then yields \(\det((Q_j)_{-x,-x})\ne0\). Thus \(\mathcal M_j=(n,Q_j,H_j,\varepsilon)\in\mathfrak F_x(P_X)\), and division of (22) by its \(x\)-coordinate proves the displayed vector response. If \(C\) is rational, Gaussian elimination can choose rational extension rows; inversion, the shear, row permutation, and division by nonzero diagonal entries preserve rationality, proving the final assertion.